\printedpage{29}
% head: A. Grothendieck : Dérivateurs, Ch. XII, Édition des Universités de Montpellier II et Paris VII
whereas all the derived notions, as well as the notions of fibrations to boot, have an unfortunate tendency \textit{not} to be stable under composition.

We could call pseudo-$W$-proper resp. quasi-$W$-proper (a little better than pseudo) what we have just brought out - I hear snickering! And if we want both at once - $\lambda$ and $\lambda^{\circ}$ associated as we said? I see the same principle of description with $S''\fl S'$ over $S$, of requiring that $X''\fl X'$ be transformed into \texttt{[elements of]} $W$. With the following choices which look nice

\begin{itemize}
\item[a$_1$)] \msp{page 39}$Y'$ has a final object, $Y''=e$. The definition thus comes down to the condition that the
\[ X_s\fl X/s, \qquad X/s'\fl X/s \]
(for $s\in\Ob S$ where $u : s'\fl s$ in $\mathrm{S}$) be in $W$, and that it stay that way after every base change $S'\fl S$.

\item[a$_2$)] $Y'$ has an initial object, $Y''=e$, this is the dual condition by dint of
\[ X_s\fl s\backslash X, \qquad s\backslash X\fl s'\backslash X. \]

\item[a$_0$)] $Y'=\boldsymbol{\Delta}^1$, $Y''=e$ (\footnote[46]{NB a$_1$), a$_2$) are denoted the same way on p. 4, but a$_0$) was forgotten.}).

\item[b)] $Y'$ is contractible, $Y''$ too (or if preferred, $Y'$ contractible, $Y''=e$)
\end{itemize}

thus the condition comes down to this: for every \textit{contractible} $Y'$ over $Y$ (notion independent of $W$, that is its main advantage) the inclusions of the fibers of $X'=X\times_Y Y'$ over $Y'$ into $X'$ are \texttt{[in]} $W$.

I feel that even for general $W$, b) remains in the realm of the verifiable, whereas we go soft if we go to the extreme of generality in this direction:

\begin{itemize}
\item[b$'$)] As b), but assuming only $Y'$ $W$-aspheric (\footnote[47]{This is essentially a$'$) of p. 4}).
\end{itemize}

Finally (condition a) of page 4)

\begin{itemize}
\item[c)] If $Y''\fl Y'$ is \texttt{[in]} $W$ without more, $X''\fl X'$ too.
\end{itemize}
% foot: Version 7 décembre 2002

\printedpage{30}
% head: A. Grothendieck : Dérivateurs, Ch. XII, Édition des Universités de Montpellier II et Paris VII
\msp{page 40}The only one of these notions that has the slightest desire to be stable under composition is the strongest c). On the other hand, the notion of a $W$-perfect arrow (cf. def. 3) is spontaneously stable under composition:

\textbf{Definition 3.} We say that $f : X\fl S$ is $W$-\textit{perfect} if it is both $W$-proper and $W$-smooth.

In the case $W=W_\infty$, all the variants on the preceding page, with the sole exception of a$_0$), are equivalent. It would be interesting to bring out minimal conditions on $W$ to ensure that this is so. (And the question arises whether in the case of $W_\infty$, a$_0$) suffices to imply the rest - I hardly expect it \dots)

In the general case, we have the diagram of implications (and even so I have limited myself)
\[
\begin{tikzcd}[column sep=small]
 & & \mathrm{c} \arrow[d, Rightarrow] & & \text{ultrastrong $W$-fibration} \\
 & & \mathrm{b}' \arrow[dd, Rightarrow] & & \text{$W$-strong $W$-fibration} \\
 & \text{$W$-perfect} \arrow[dr, Rightarrow] & & & \\
 & & \mathrm{b} \arrow[dl, Rightarrow] \arrow[dr, Rightarrow] & & \text{hot-strong $W$-fibration (\footnotemark[48])} \\
\text{left $W$-fibration} & \mathrm{a}_1 \arrow[dr, Rightarrow] & & \mathrm{a}_2 \arrow[dl, Rightarrow] & \text{right $W$-fibration} \\
 & & \mathrm{a}_0 & & \text{$W$-pseudofibration}
\end{tikzcd}
\]
\footnotetext[48]{It would seem that this is, in the end, the ``right notion'' of $W$-fibration, \textit{equivalent} to a$_1$ and a$_2$, which would be equivalent to each other. Cf. prop. 9, p. 48.}

\texttt{[``$W$-perfect $\Longrightarrow$ b'' circled]}.

\msp{page 41}\textbf{Proposition 8.} \textit{Let $f : X\fl S$ be in} Cat. \textit{The following conditions are equivalent.}

\begin{itemize}
\item[\textit{a)}] \textit{$f$ is a left $W$-fibration, i.e. condition} a$_1$ \textit{of page 39 is satisfied.}

\item[\textit{b)}] \textit{$f$ is $W$-preproper, and $\forall\; s_0\fl s_1$ in $S$, $X/s_0\fl X/s_1$ is in $W$.}

\item[\textit{c)}] \textit{$f$ is $W$-preproper, and a $W$-pseudofibration, i.e. condition a$_0$ of page 39 is satisfied.} % AUDIT resolved: CMTI12 text-italic a with CMR8 subscript; see corrections.csv, PDF leaf 40, ms41
\end{itemize}

(\footnote[49]{\textbf{Corollary.} \textit{When $f$ is} Cat-\textit{cofibrant, then these conditions are also equivalent to this: $\forall\; u : s_0\fl s_1$ in $S$, $u_* : X_{s_0}\fl X_{s_1}$ is in $W$.}}).
% foot: Version 7 décembre 2002

\printedpage{31}
% head: A. Grothendieck : Dérivateurs, Ch. XII, Édition des Universités de Montpellier II et Paris VII
\textbf{NB} These conditions are satisfied if $f$ is at the same time $W$-presmooth, cf. prop. 7.

\textsc{Proof}. We have a) $\Longrightarrow$ b) trivially (take $Y''=Y/s$, $Y''=Y/s'$ \texttt{[connection with the assertion? Cf. p. 4.]}).

b) $\Longleftrightarrow$ c) by the reasoning of proposition 7, for in terms of $X'=X_u=X\times_S(\boldsymbol{\Delta}^1,u)$, condition b) for a given $u : s_0\fl s_1$ is interpreted in the commutative diagram
\[
\begin{tikzcd}[column sep=small, row sep=small]
X'_0 \arrow[dr, "i'_0\in W"] \arrow[dd, "\wr"] & & X'_1 \arrow[dr, "i'_1\in W"] \arrow[dd, "\wr"] & \\
 & \underset{(\simeq X'_0)}{X'/0} \arrow[rr, "\varphi'", crossing over] \arrow[dd, "p_0"] & & X'/1\ (\simeq X') \arrow[dd, "p_1"] \\
X_{s_0} \arrow[dr, "i_0\in W"'] & & X_{s_1} \arrow[dr, "i_1\in W"] & \\
 & X/s_0 \arrow[rr, "\varphi"] & & X/s_1
\end{tikzcd}
\]

\texttt{[there are a few redundant $W$ in the original diagram]} where $i_0$, $i_1$, $i'_0$, $i'_1$ are in $W$ because of the hypothesis of $W$-preproperness (stable under base change).

Thus $p_0$, $p_1$ are \texttt{[in]} $W$, thus $\varphi\in W$ if and only if $\varphi'\in W$, now $\varphi'\in W$ is condition a$_0$.

\msp{page 42}Let us finally prove that b), c) imply a). Since c) (hence also b)) is stable under base change, it suffices to show this: If $S$ has a final object $e$, then for every $s\in\Ob S$, the inclusion $X_s\xrightarrow{\ i\ }X$ is in $W$. Now we have the commutative diagram
\[
\begin{tikzcd}
X_s \arrow[r, "i_s\in W"] \arrow[dr, "i"'] & X/s \arrow[d, "u\,\in\,W\text{ by c) }\texttt{[rather by b)]}"] \\
 & X/e=X,
\end{tikzcd}
\]
which indeed shows $i=ui_s\in W$.

Dually, if $f$ is $W$-presmooth, then condition a$_0$ ($f$ a $W$-pseudofibration) is equivalent to a$_2$ ($f$ a right $W$-fibration), and also means that the
\[ s_1\backslash X\fl s_0\backslash X \]
are in $W$. But the notion of ``$W$-fibration'' that we land on (if we are not content with the paltry a$_0$!) is different in the two cases dual to each other. We would like in both cases to find ``hot-strong'' $W$-fibrations (condition b), at least by going all out (if necessary) - i.e. for $f$ $W$-perfect. So it comes down to this.

\textbf{Question:} \textit{Let $X\xrightarrow{\ f\ }S$, with $f$ $W$-perfect, $S$ contractible, $s\in\Ob S$. The inclusion $X_s\fl X$, is it in $W$?}
% foot: Version 7 décembre 2002

\printedpage{32}
% head: A. Grothendieck : Dérivateurs, Ch. XII, Édition des Universités de Montpellier II et Paris VII
\msp{page 43}We are going to assume only that $f$ is $W$-smooth and a left $W$-fibration \texttt{[rather that $f$ is simply a right $W$-fibration, i.e that $f$ is $W$-presmooth and a $W$-pseudofibration]} (cf. prop. 8). We are going to reduce to the case where $X$ is even fibrant over $S$, by introducing $\Phi_S(X)$ \texttt{[fibered category defined by the functor $S^{\circ}\fl\Cat$, $s\longmapsto s\backslash X$]} and % [sic] i.e que -> i.e. que, missing period of the abbreviation
\[ X\xrightarrow{\ i_X\ }\Phi_S(X). \]

The trouble is that I have done the theory of $\HOT_W$ up to now with the hypothesis (which turns out to be idiotic) that $W\supseteq W_\infty$ - I will therefore be obliged to review quickly under what more reasonable hypotheses on $W$ the familiar results are valid. For example I would like $i_X\in W$ (\footnote[50]{in fact, $i_X\in(W_\omega)^{\mathrm{d}}_S\subseteq W^{\mathrm{d}}_S\subseteq W_S$.}), so that we have
\[
\begin{tikzcd}[column sep=large, row sep=large]
X_s \arrow[r, "i_s\in W", "\text{since $X$ $W$-smooth}"'] \arrow[d, "j^X_s"'] & \Phi_S(X)_s=s\backslash X \arrow[d, "j_s^{\Phi_S(X)}"] \\
X \arrow[r, "i_X\in W"] & \Phi_S(X),
\end{tikzcd}
\]
we therefore see that to prove that the inclusion $j^X_s$ is \texttt{[in]} $W$, it suffices to prove the same thing for $\Phi_S(X)$, which is fibrant (not just $W$-smooth). Moreover, the pseudofibration hypothesis on $X$ surely passes to $\Phi_S(X)$ (I leave the two verifications for later (\footnote[51]{O.K.})). We are thus reduced to the case of
\[ X\in\Ob\uFib_0S, \qquad S\text{ contractible,} \]
to prove that the inclusion
\[ X_s\fl X \]
\msp{page 44}is in $W$. I am going to prove this:

\textbf{Lemma.} \textit{Let $S$ be in} Cat, \textit{$X\in\Ob\uFib_0S$, and let us consider, for an $S'$ in} Cat,
\[
\begin{aligned}
f &\quad\longmapsto\quad f^*(X/S)=X\times_S(S',f)\\
\uHom(S',S)^{\circ} &\quad\longrightarrow\quad \uFib_0S'.
\end{aligned}
\]
\textit{(This functor is defined thanks to the fact that $X$ is fibrant over $S$, and can therefore be interpreted, up to $S$-equivalence, in terms of a functor $S^{\circ}\fl$} Cat \textit{\dots\ (\footnote[52]{\textbf{NB} $\Phi_S(X)$ was even already split, defined by the functor $s\longmapsto s\backslash X$.})). The arrows of $\uHom(S',S)$ are transformed into arrows $\in W^{\mathrm{d}}_{S'}$.}
% foot: Version 7 décembre 2002

\printedpage{33}
% head: A. Grothendieck : Dérivateurs, Ch. XII, Édition des Universités de Montpellier II et Paris VII
Indeed, in terms of functors
\[ S'^{\circ}\fl\Cat, \]
if $X$ is defined by
\[ S^{\circ}\xrightarrow{\ \mathscr{X}\ }\Cat, \]
then $f^*(X)$ is defined by $\mathscr{X}\circ f$, and the arrow
\[ \mathscr{X}\circ f\fl\mathscr{X}\circ g \]
defined by $u : f\fl g$ is none other than $\mathscr{X}*u$. Its value at $s'\in\Ob S'$ is obtained by taking
\[ u(s') : f(s')\fl g(s'), \]
whence by $\mathscr{X}$
\[ \mathscr{X}(u(s')) : \underbrace{\mathscr{X}(g(s'))}_{(\mathscr{X}\circ g)(s')}\fl\underbrace{\mathscr{X}(f(s'))}_{(\mathscr{X}\circ f)(s')} \]
(\footnote[53]{to bring out a separate statement.}). But the hypothesis $X\in\uFib_0S$ means that the $\mathscr{X}(v)$, $v\in\mathrm{Fl}(S)$, are all in $W$. Thus
\msp{page 45}$\mathscr{X}$ is in $W$ fiber by fiber, thus in $W^{\mathrm{d}}_{S'}$.

\textbf{Corollary 1.} \textit{The functor of the lemma, composed with}
\[ \uFib_0S'\fl\HOT^{\mathrm{lc}}_W(S')\overset{\text{def}}{=}(\uFib_0S')(W^{\mathrm{d}}_{S'})^{-1} \]
\textit{factors through the fundamental groupoid of} $\uHom$ \textit{:}
\[
\begin{tikzcd}[row sep=huge]
\uHom(S',S)^{\circ} \arrow[r] \arrow[d] & \uFib_0S' \arrow[d] \\
\textstyle\prod(\uHom(S',S))^{\circ} \arrow[r] & \HOT^{\mathrm{lc}}_W(S').
\end{tikzcd}
\]
\textit{If $f$, $g$ are in the same connected component of $\uHom(S',S)$, then $f^*X$ and $g^*X$ are isomorphic (noncanonically, in general) in} $\HOT^{\mathrm{lc}}_W(S')$.

It is cleverer not to fix $X$ and to consider
\[ \uHom(S',S)\fl\uHom(\HOT^{\mathrm{lc}}_W(S),\HOT^{\mathrm{lc}}_W(S')) \tag*{$(*)$} \]
corresponding to
\[
\begin{aligned}
\uHom(S',S) &\quad\times\quad \HOT^{\mathrm{lc}}_W(S) &&\quad\longrightarrow\quad &&\HOT^{\mathrm{lc}}_W(S')\\
(f &\quad,\quad\ \xi) &&\quad\longmapsto\quad && f^*(\xi).
\end{aligned}
\]
The lemma will say that the arrows of $S'\fl S$ are transformed by $(*)$ into isomorphisms, whence
\[ \prod\uHom(S',S)\fl\uHom(\HOT^{\mathrm{lc}}_W(S),\HOT^{\mathrm{lc}}_W(S')), \]
% foot: Version 7 décembre 2002